% ======================================================================
% SECTION 6: CONCLUSIONS
% Summary of accomplishments and persistent themes.
% ======================================================================

This paper documented four author Physical AI oncology clinical trial
simulations and showed how their cumulative artifact set demonstrates
a computational and contextual capability that distinguishes LLM-
scale repository-context prediction from the narrow supervised
oncology models named in \S\ref{sec:intro-ai-baseline}.

\textbf{Headline artifact counts.} Simulation 1 produced 56 hours of
continuous trial coverage (\path{hour-00} through \path{hour-55})
with 7 files per hour for a total of 392 site-side artifacts at
minute resolution across 4 sites and 116 robots. Simulation 2
produced 10 stage Python scripts plus 30 ASCII progress diagrams plus
6 deliverable diagrams plus 6 regulatory tables for one patient
(PAT-2026-0042) over 1,120 days. Simulation 3 produced 24 hourly
Python scripts, 24 JSON outputs, 75 ASCII diagrams, and 53 core
agents organized into 4 layers. Simulation 4 produced 168 hourly
Python scripts, 168 JSON outputs, 525 ASCII diagrams, 7 daily
summaries, 7 GitHub branches, 168 GitHub commits, 7 pull requests,
plus a complete local-verification run on a 2015 Intel Core i5-6200U
laptop with 4 GB of RAM running Windows 10 Pro. Simulations 1 and 2
are clinical trial site simulations; Simulations 3 and 4 are clinical
trial sponsor simulations. The site and sponsor simulations are
synergistic: a sponsor simulation can ingest the signal stream from a
site simulation, and a site simulation can ingest the authorization
stream from a sponsor simulation, with the FDA RTCT real-time channel
observing both streams.

\textbf{Persistent themes.} Three themes recur across the four
simulations. The first is repository-scale LLM context (1M tokens) as
the computational signature that distinguishes this work from current
narrow supervised oncology models; the model holds protocol,
regulatory, robotic, EHR, and ASCII-diagram content in a single
inference pass and emits multiple artifact types in the same step.
The second is hourly commit cadence as the operational signature
aligned with the FDA RTCT vision; each simulation emits at least one
commit per hour to a public GitHub repository, producing an audit
trail that mirrors the real-time signal-sharing framework the FDA
validated through Paradigm Health \cite{fda2026realtime}. The third
is the file-level artifact bundle (ASCII plus Markdown plus Python
plus JSON) that makes the work auditable and human-monitorable in
real time. A clinical reader can read the ASCII facility diagram
without running any code; a downstream analytics agent can import
the Python module directly; an FDA inspector can verify the GitHub
commit log against the local re-run.

\textbf{Implications for patient safety and effectiveness.} LLM
code-and-text contextual awareness at scale enables the strongest
opportunity to improve patient prediction safety and effectiveness in
oncology trials, because (a) the model can reason across protocol,
EHR, robotic telemetry, regulatory adaptations, and ASCII diagrams in
a single inference pass; (b) the model can emit structured artifacts
(JSON, Python, ASCII) that downstream systems can ingest; and (c) the
model's output can be committed to GitHub at a 1-hour cadence,
mirroring the real-time signal-sharing framework the FDA validated
through Paradigm Health \cite{fda2026realtime}. The framing here is
respectful of the FDA and of current trial sponsors; the four
simulations are opportunistic extensions of the FDA RTCT vision, not
critiques of it. The advanced robotics integration (116 robot
instances in Simulation 1) and the advanced predictive layer (1M
token context across all four simulations) extend the FDA's
pharmacology-only pair of proof-of-concept trials toward a more
complete continuous-trial framework, which is crucial for both
patient safety (real-time AE detection across robotics and pharma)
and patient efficacy (per-hour outcome forecasting at the protocol
level).

\textbf{Key limitations restated.} None of the four simulations was
tested on real patients. Claude Code Opus 4.7 Max output is non-
deterministic. The local verification on the Core i5-6200U laptop
was partial. The four simulations do not satisfy TRIPOD+AI or CREMLS
reporting requirements \cite{Collins2024TRIPODAI, ElEmam2024CREMLS}.
The work is a computational and software demonstration, not a
clinical demonstration.

\textbf{Forward path.} The next step is to bridge the synthetic-to-
real gap. The Track A and Track B future-work directions in
\S\ref{sec:limits-future} provide the architectural plan for a real-
patient extension of this work, with three concrete deliverables: a
TRIPOD+AI-compliant retrospective validation of one of the four
simulations on a real-patient cohort, a public benchmark against the
supervised baselines named in Background-A and Background-B, and an
FDA-aligned RTCT pilot submission that uses Simulation 4 as the
sponsor-side input. The artifact at
\path{new-trial/national-24-7-trial/paper/full-paper/LaTeX_Source_Files.zip}
is ready for upload to Overleaf, and the GitHub commit log of this
pull request provides the audit trail for the present
template-population pass. Cite \cite{fda2026realtime,
kawchak_2026_19176370, kawchak_2026_19244918, kawchak_2026_18810541,
kawchak_2026_19119939, kawchak_2026_19396256, Collins2024TRIPODAI,
ElEmam2024CREMLS} for the artifacts and reporting standards that
underpin this conclusion.
